% GENERATED FILE. Edit canonical lessons, then run npm run generate:books.
% canonical-source-hash: fnv1a64:49f9116e4be4d2cc
% canonical-lessons: ZH-C84-yiqi, ZH-C84-jianmian, ZH-C84-liaotian, ZH-C84-canjia, ZH-C84-zhunbei, ZH-R84-review-80-84, ZH-R84-review-final-84

\chapter{Doing Things Together}
\label{ch:zh-doing-things-together}
\hlchaptermodality{84}

\begin{chapteropening}
\textbf{By the end of this chapter:} \emph{I can arrange to meet, chat, take part and get ready together.}

Complete the last lesson of chapter 84: I can arrange to meet, chat, take part and get ready together.
\end{chapteropening}

\section[\texorpdfstring{\zh{一起} (yìqǐ)}{一起 (yìqǐ)}]{\zh{一起} (yìqǐ) — together}

\label{lesson:ZH-C84-yiqi}



\begin{quote}
Before the new one: say the Mandarin for to run out of time, then the Mandarin for to be late.
\end{quote}

\subsection*{You'll want to know: \zh{一起}}
\textbf{\zh{一起}} — \emph{yìqǐ} — \textquotedblleft{}together\textquotedblright{}.

\begin{quote}
\textbf{\zh{三年前我们一起上学}} \emph{sān nián qián wǒmen yìqǐ shàngxué} — Three years ago we were at school together.
\end{quote}

yī changes its tone here, falling to yì before the dipping qǐ: yìqǐ.

\begin{cousinweb}
\textbf{\zh{一起}} (\emph{yìqǐ}) is built from \zh{一} 'one' and \zh{起} 'to rise', both inherited from Middle Chinese, \emph{'jit} and \emph{khiX} (Guangyun, in Baxter's transcription). Rising as one is a likely picture, but how the pair came to mean together is not settled.
\end{cousinweb}

\begin{grammarlens}[title={What you've built}]
One more word for everyday answers.
\end{grammarlens}

\subsection*{Your turn}
\begin{itemize}
\raggedright
  \item \emph{Say it:} \emph{yìqǐ}
  \item \emph{Say it:} \emph{yìqǐ}, once more
  \item \emph{Say it:} \emph{sān nián qián wǒmen yìqǐ shàngxué} — Three years ago we were at school together.
  \item \emph{Recall:} say the Mandarin for to run out of time, then the Mandarin for to be late, then say \emph{yìqǐ} again
\end{itemize}

\begin{tcolorbox}[breakable,colback=teal!4,colframe=teal!35!black,title={Before you move on}]
What does \zh{一起} mean? (\textquotedblleft{}Together\textquotedblright{}.) Say it once more.
\end{tcolorbox}
\section[jiànmiàn]{jiànmiàn — to meet up; to see each other}

\label{lesson:ZH-C84-jianmian}



\begin{quote}
Before the new one: say the Mandarin for to be late, then the Mandarin for together.
\end{quote}

\subsection*{You'll want to know: jiànmiàn}
\emph{jiànmiàn} — \textquotedblleft{}to meet up; to see each other\textquotedblright{}.

\begin{quote}
\emph{wǒmen xīngqītiān jiànmiàn} — We're meeting on Sunday.
\end{quote}

Its characters come later; for now it is a word for the ear.

\begin{cousinweb}
\emph{jiànmiàn} is built from \zh{见} 'to see' and miàn 'face', both inherited from Middle Chinese, \emph{kenH} and \emph{mjienH} (Guangyun, in Baxter's transcription): to see each other's face.
\end{cousinweb}

\begin{grammarlens}[title={What you've built}]
One more word for everyday answers.
\end{grammarlens}

\subsection*{Your turn}
\begin{itemize}
\raggedright
  \item \emph{Say it:} \emph{jiànmiàn}
  \item \emph{Say it:} \emph{jiànmiàn}, once more
  \item \emph{Say it:} \emph{wǒmen xīngqītiān jiànmiàn} — We're meeting on Sunday.
  \item \emph{Recall:} say the Mandarin for to be late, then the Mandarin for together, then say \emph{jiànmiàn} again
\end{itemize}

\begin{tcolorbox}[breakable,colback=teal!4,colframe=teal!35!black,title={Before you move on}]
What does jiànmiàn mean? (\textquotedblleft{}To meet up; to see each other\textquotedblright{}.) Say it once more.
\end{tcolorbox}
\section[liáotiān]{liáotiān — to chat}

\label{lesson:ZH-C84-liaotian}



\begin{quote}
Before the new one: say the Mandarin for together, then the Mandarin for to meet up.
\end{quote}

\subsection*{You'll want to know: liáotiān}
\emph{liáotiān} — \textquotedblleft{}to chat\textquotedblright{}.

\begin{quote}
\emph{wǒmen hē chá liáotiān} — We drink tea and chat.
\end{quote}

Its characters come later; for now it is a word for the ear.

\begin{cousinweb}
\emph{liáotiān} is built from liáo and \zh{天} 'sky', both inherited from Middle Chinese, \emph{lew} and \emph{then} (Guangyun, in Baxter's transcription). The dictionary glosses the first as a helping word; how the pair came to mean chatting is not settled.
\end{cousinweb}

\begin{grammarlens}[title={What you've built}]
One more word for everyday answers.
\end{grammarlens}

\subsection*{Your turn}
\begin{itemize}
\raggedright
  \item \emph{Say it:} \emph{liáotiān}
  \item \emph{Say it:} \emph{liáotiān}, once more
  \item \emph{Say it:} \emph{wǒmen hē chá liáotiān} — We drink tea and chat.
  \item \emph{Recall:} say the Mandarin for together, then the Mandarin for to meet up, then say \emph{liáotiān} again
\end{itemize}

\begin{tcolorbox}[breakable,colback=teal!4,colframe=teal!35!black,title={Before you move on}]
What does liáotiān mean? (\textquotedblleft{}To chat\textquotedblright{}.) Say it once more.
\end{tcolorbox}
\section[cānjiā]{cānjiā — to take part in; to join}

\label{lesson:ZH-C84-canjia}



\begin{quote}
Before the new one: say the Mandarin for to meet up, then the Mandarin for to chat.
\end{quote}

\subsection*{You'll want to know: cānjiā}
\emph{cānjiā} — \textquotedblleft{}to take part in; to join\textquotedblright{}.

\begin{quote}
\emph{nǐ xiǎng cānjiā ma} — Would you like to join in?
\end{quote}

Its characters come later; for now it is a word for the ear.

\begin{cousinweb}
\emph{cānjiā} is built from cān 'to attend' and jiā 'to add', both inherited from Middle Chinese, \emph{tshom} and \emph{kae} (Guangyun, in Baxter's transcription).
\end{cousinweb}

\begin{grammarlens}[title={What you've built}]
One more word for everyday answers.
\end{grammarlens}

\subsection*{Your turn}
\begin{itemize}
\raggedright
  \item \emph{Say it:} \emph{cānjiā}
  \item \emph{Say it:} \emph{cānjiā}, once more
  \item \emph{Say it:} \emph{nǐ xiǎng cānjiā ma} — Would you like to join in?
  \item \emph{Recall:} say the Mandarin for to meet up, then the Mandarin for to chat, then say \emph{cānjiā} again
\end{itemize}

\begin{tcolorbox}[breakable,colback=teal!4,colframe=teal!35!black,title={Before you move on}]
What does cānjiā mean? (\textquotedblleft{}To take part in; to join\textquotedblright{}.) Say it once more.
\end{tcolorbox}
\section[zhǔnbèi]{zhǔnbèi — to get ready; to prepare}

\label{lesson:ZH-C84-zhunbei}



\begin{quote}
Before the new one: say the Mandarin for to chat, then the Mandarin for to take part.
\end{quote}

\subsection*{You'll want to know: zhǔnbèi}
\emph{zhǔnbèi} — \textquotedblleft{}to get ready; to prepare\textquotedblright{}.

\begin{quote}
\emph{wǒ zhǔnbèi chànggē} — I'm getting ready to sing.
\end{quote}

Its characters come later; for now it is a word for the ear.

\begin{cousinweb}
\emph{zhǔnbèi} is built from zhǔn 'level, a standard' and bèi 'to have ready', both inherited from Middle Chinese, \emph{tsywinX} and \emph{bijH} (Guangyun, in Baxter's transcription).
\end{cousinweb}

\begin{grammarlens}[title={What you've built}]
One more word for everyday answers.
\end{grammarlens}

\subsection*{Your turn}
\begin{itemize}
\raggedright
  \item \emph{Say it:} \emph{zhǔnbèi}
  \item \emph{Say it:} \emph{zhǔnbèi}, once more
  \item \emph{Say it:} \emph{wǒ zhǔnbèi chànggē} — I'm getting ready to sing.
  \item \emph{Recall:} say the Mandarin for to chat, then the Mandarin for to take part, then say \emph{zhǔnbèi} again
\end{itemize}

\begin{tcolorbox}[breakable,colback=teal!4,colframe=teal!35!black,title={Before you move on}]
What does zhǔnbèi mean? (\textquotedblleft{}To get ready; to prepare\textquotedblright{}.) Say it once more.
\end{tcolorbox}
\section[(dialogue)]{Looking back: from Jobs to Doing Things Together}

\label{lesson:ZH-R84-review-80-84}



\begin{quote}
Say the Mandarin for to go to work and worker.
\end{quote}

\subsection*{Your turn}
\begin{itemize}
\raggedright
  \item \emph{Say it:} to go to work, worker, driver, doctor, company
  \item \emph{Say it:} inside, city, place, abroad, beside
  \item \emph{Say it:} after, the left side, the right side, the middle, opposite
  \item \emph{Say it:} time to spare, early, right away, to run out of time, to be late
  \item \emph{Say it:} together, to meet up, to chat, to take part, to get ready
  \item \emph{Read it:} \zh{里} — each one aloud, then a word that holds it
\end{itemize}

\begin{tcolorbox}[breakable,colback=teal!4,colframe=teal!35!black,title={Before you move on}]
To go to work? (\emph{shàngbān}.) Worker? (\emph{gōngrén}.) Driver? (\emph{sījī}.) Doctor? (\emph{yīshēng}.) Company? (\emph{gōngsī}.) Inside? (\textbf{\zh{里}}, \emph{lǐ}.) City? (\emph{chéngshì}.) Place? (\emph{dìfang}.) Abroad? (\emph{wàiguó}.) Beside? (\emph{pángbiān}.) After? (\emph{hòu}.) The left side? (\emph{zuǒbian}.) The right side? (\emph{yòubian}.) The middle? (\emph{zhōngjiān}.) Opposite? (\emph{duìmiàn}.) Time to spare? (\emph{shíjiān}.) Early? (\textbf{\zh{早}}, \emph{zǎo}.) Right away? (\textbf{\zh{马上}}, \emph{mǎshàng}.) To run out of time? (\emph{láibují}.) To be late? (\emph{chídào}.) Together? (\textbf{\zh{一起}}, \emph{yìqǐ}.) To meet up? (\emph{jiànmiàn}.) To chat? (\emph{liáotiān}.) To take part? (\emph{cānjiā}.) To get ready? (\emph{zhǔnbèi}.)
\end{tcolorbox}
\section[(dialogue)]{Doing Things Together — once more}

\label{lesson:ZH-R84-review-final-84}



\begin{quote}
Say the Mandarin for together and to meet up.
\end{quote}

\subsection*{Your turn}
\begin{itemize}
\raggedright
  \item \emph{Say it:} together, to meet up, to chat, to take part, to get ready
\end{itemize}

\begin{tcolorbox}[breakable,colback=teal!4,colframe=teal!35!black,title={Before you move on}]
Together? (\textbf{\zh{一起}}, \emph{yìqǐ}.) To meet up? (\emph{jiànmiàn}.) To chat? (\emph{liáotiān}.) To take part? (\emph{cānjiā}.) To get ready? (\emph{zhǔnbèi}.)
\end{tcolorbox}
